\section*{Bab \ref{ch:g12:proton-nmr} --- NMR Proton}

\begin{solution}{exo:g12:proton-nmr:1}
Metana: 1. Etana: 1 (kedua \ce{CH3} serupa). Propana: 2 (dua \ce{CH3}
yang serupa, satu \ce{CH2}). Metanol: 2 (\ce{CH3} dan \ce{OH}).
Metoksimetana: 1.
\end{solution}

\begin{solution}{exo:g12:proton-nmr:2}
2 tetangga: triplet, 1:2:1. 3 tetangga: kuartet, 1:3:3:1. 0 tetangga:
singlet.
\end{solution}

\begin{solution}{exo:g12:proton-nmr:3}
Total \qty{40}{mm} untuk 8 \omterm{def:g9:inside-the-atom:proton}{proton}: \qty{5}{mm} per \omterm{def:g9:inside-the-atom:proton}{proton}. Sinyalnya
mewakili 3, 2, dan 3 \omterm{def:g9:inside-the-atom:proton}{proton}.
\end{solution}

\begin{solution}{exo:g12:proton-nmr:4}
\omterm{def:g11:functional-groups:carbonyl}{Aldehida}: 9.7 sampai \qty{10.0}{ppm}. \ce{CH3} dalam rantai: 0.7 sampai
\qty{1.3}{ppm}.
\end{solution}

\begin{solution}{exo:g12:proton-nmr:5}
Kedua belas \omterm{def:g12:proton-nmr:equivalent}{proton ekuivalennya} memberikan satu garis tunggal yang
tajam dan kuat, mudah ditempatkan di nol. \omterm{def:g6:solutions-and-solubility:solute}{Pelarut} terdeuterasi tidak
memberikan sinyal \omterm{def:g9:inside-the-atom:proton}{proton} sendiri, yang kalau tidak demikian akan
menenggelamkan sinyal beberapa miligram sampel itu.
\end{solution}

\begin{solution}{exo:g12:proton-nmr:6}
Dua sinyal. \ce{CH2} di sebelah klorin: 2 \omterm{def:g9:inside-the-atom:proton}{proton}, kuartet (3 tetangga),
dalam rentang 2.5--\qty{4.0}{ppm}. \ce{CH3}: 3 \omterm{def:g9:inside-the-atom:proton}{proton}, triplet (2
tetangga), dekat rentang \ce{CH3} yang biasa, sedikit lebih tinggi
karena klorinnya berjarak dua \omterm{def:g7:atoms-and-molecules:atom}{atom}.
\end{solution}

\begin{solution}{exo:g12:proton-nmr:7}
Tiga sinyal: kedua \ce{CH3}, 6 \omterm{def:g12:proton-nmr:equivalent}{proton ekuivalen}, doblet (1 tetangga,
yaitu \ce{CH}); \ce{CH}, 1 \omterm{def:g9:inside-the-atom:proton}{proton}, dipecah oleh 6 tetangga menjadi 7
garis, dalam rentang \ce{H-C-O} (3.3--\qty{4.5}{ppm}); \ce{OH}, 1
\omterm{def:g9:inside-the-atom:proton}{proton}, singlet.
\end{solution}

\begin{solution}{exo:g12:proton-nmr:8}
(a) propanon; (b) etil etanoat; (c) asam propanoat.
\end{solution}

\begin{solution}{exo:g12:proton-nmr:9}
\ce{CH2} terikat pada \omterm{def:g7:atoms-and-molecules:atom}{atom} oksigen \omterm{def:g11:functional-groups:carboxylic-acid}{ester}, yang elektronegatif dan
mengurangi perlindungan \omterm{def:g9:inside-the-atom:proton}{proton-protonnya}: rentang \ce{H-C-O},
3.3--\qty{4.5}{ppm}. \ce{CH3} berjarak satu karbon lebih jauh dan tetap
dalam rentang rantai.
\end{solution}

\begin{solution}{exo:g12:proton-nmr:10}
Propanon: sebuah \ce{C=O} (\omterm{def:g11:functional-groups:carbonyl}{keton}, \qty{1715}{cm^{-1}}) dan enam \omterm{def:g12:proton-nmr:equivalent}{proton
ekuivalen}, satu singlet. Propanal akan memberikan tiga sinyal, salah
satunya, \omterm{def:g9:inside-the-atom:proton}{proton} \omterm{def:g11:functional-groups:carbonyl}{aldehida} \ce{CH3-CH2-CHO}, di dekat
9.7--\qty{10.0}{ppm}.
\end{solution}

\begin{solution}{exo:g12:proton-nmr:11}
\omterm{def:g9:inside-the-atom:proton}{Proton} \ce{OH} dipertukarkan dengan cepat antarmolekul: \omterm{def:g9:inside-the-atom:proton}{proton} itu
memberikan singlet dan tidak memecah sinyal tetangganya. Karena itu
\ce{CH2} hanya dipecah oleh 3 \omterm{def:g9:inside-the-atom:proton}{proton} \ce{CH3}: kuartet.
\end{solution}

\begin{solution}{exo:g12:proton-nmr:12}
Asam butanoat \ce{CH3-CH2-CH2-COOH}: 4 sinyal, triplet (\ce{CH3}),
sinyal 6 garis (\ce{CH2} tengah, 5 tetangga), triplet (\ce{CH2-C=O}),
singlet di dekat 11--\qty{12}{ppm}. Etil etanoat: singlet, kuartet,
triplet. Metil propanoat \ce{CH3-CH2-CO-O-CH3}: triplet, kuartet
(\ce{CH2-C=O}, 2.0--\qty{2.4}{ppm}), singlet (\ce{O-CH3}). Propil
metanoat \ce{H-CO-O-CH2-CH2-CH3}: 4 sinyal, singlet jauh di kiri (H
pada karbon \ce{C=O}), triplet (\ce{O-CH2}), sinyal 6 garis, triplet
(\ce{CH3}).
\end{solution}

\begin{solution}{exo:g12:proton-nmr:13}
Sinyal \ce{OH}: \omterm{def:g9:inside-the-atom:proton}{proton-proton} \ce{O-H} bertukar dengan deuterium air
berat, dan \ce{O-D} tidak memberikan sinyal \omterm{def:g9:inside-the-atom:proton}{proton}. Sinyal yang hilang
setelah dikocok dengan \ce{D2O} adalah milik \omterm{def:g9:inside-the-atom:proton}{proton} pada O atau N.
\end{solution}

\begin{solution}{exo:g12:proton-nmr:14}
Etanol menambahkan kuartet di dekat \qty{3.69}{ppm}, triplet di dekat
\qty{1.23}{ppm}, dan singlet \ce{OH}. Singlet \omterm{def:g11:functional-groups:carboxylic-acid}{ester}: \qty{30}{mm} untuk
3 \omterm{def:g9:inside-the-atom:proton}{proton}, \qty{10}{mm} per \omterm{def:g9:inside-the-atom:proton}{proton} \omterm{def:g11:functional-groups:carboxylic-acid}{ester}. Triplet etanol: \qty{3}{mm}
untuk 3 \omterm{def:g9:inside-the-atom:proton}{proton}, \qty{1}{mm} per \omterm{def:g9:inside-the-atom:proton}{proton} etanol. Karena setiap \omterm{def:g7:atoms-and-molecules:molecule}{molekul}
menyumbangkan satu \ce{CH3} pada sinyal-sinyal ini, perbandingan
jumlahnya $1/10$: sekitar satu \omterm{def:g7:atoms-and-molecules:molecule}{molekul} etanol untuk sepuluh \omterm{def:g7:atoms-and-molecules:molecule}{molekul}
\omterm{def:g11:functional-groups:carboxylic-acid}{ester}.
\end{solution}

\begin{solution}{exo:g12:proton-nmr:15}
Empat sinyal: kedua \ce{CH3} (6 \omterm{def:g9:inside-the-atom:proton}{proton}), doblet; \ce{CH} (1 \omterm{def:g9:inside-the-atom:proton}{proton}),
dipecah oleh 6 + 2 = 8 tetangga menjadi 9 garis; \ce{CH2} (2 \omterm{def:g9:inside-the-atom:proton}{proton}) di
sebelah \ce{O}, doblet, dalam rentang \ce{H-C-O}; \ce{OH} (1 \omterm{def:g9:inside-the-atom:proton}{proton}),
singlet. Sinyal \ce{CH} memiliki garis paling banyak.
\end{solution}

\begin{solution}{pb:g12:proton-nmr:1}
\textbf{1.} \omterm{def:g11:organic-skeletons:alkane}{Alkana} dengan empat karbon adalah \ce{C4H10}; satu \ce{C=O}
menghilangkan dua hidrogen dan kedua oksigen tidak menambah apa-apa:
\ce{C4H8O2}, rumus asam dan \omterm{def:g11:functional-groups:carboxylic-acid}{ester} dengan empat karbon.

\textbf{2.} Asam butanoat \ce{CH3-CH2-CH2-COOH}; etil etanoat
\ce{CH3-CO-O-CH2-CH3}; metil propanoat \ce{CH3-CH2-CO-O-CH3}; propil
metanoat \ce{H-CO-O-CH2-CH2-CH3}.

\textbf{3.} Ketiga \omterm{def:g11:functional-groups:carboxylic-acid}{ester} memiliki gugus \omterm{def:g11:functional-groups:carboxylic-acid}{ester} yang sama; keempatnya
memiliki sebuah \ce{C=O}.

\textbf{4.} Sebuah \ce{C=O}, dalam rentang \omterm{def:g11:functional-groups:carboxylic-acid}{ester} (di dekat
\qty{1735}{cm^{-1}}).

\textbf{5.} Pita \ce{O-H} yang sangat lebar dari 2500 sampai
\qty{3300}{cm^{-1}}: tidak ada, jadi bukan asam butanoat.

\textbf{6.} Tidak: ketiga \omterm{def:g11:functional-groups:carboxylic-acid}{ester} itu memiliki gugus-gugus yang sama.

\textbf{7.} Tiga.

\textbf{8.} Sebuah gugus etil \ce{CH3-CH2-}: \ce{CH2} dipecah oleh 3
\omterm{def:g9:inside-the-atom:proton}{proton} (kuartet), \ce{CH3} oleh 2 (triplet).

\textbf{9.} \omterm{def:g9:inside-the-atom:proton}{Proton-protonnya} tidak memiliki \omterm{def:g9:inside-the-atom:proton}{proton} pada \omterm{def:g7:atoms-and-molecules:atom}{atom-atom}
tetangganya: sebuah \ce{CH3} yang terikat pada karbon \ce{C=O} atau
pada oksigen \omterm{def:g11:functional-groups:carboxylic-acid}{ester}.

\textbf{10.} \ce{CH2} itu terikat pada oksigen (rentang \ce{H-C-O}).

\textbf{11.} Triplet (\ce{CH3}), kuartet di dekat 2.0--\qty{2.4}{ppm}
(\ce{CH2} di sebelah \ce{C=O}), dan singlet dalam rentang \ce{H-C-O}
(\ce{O-CH3}). Kuartetnya akan jauh di bawah \qty{4.12}{ppm} dan
singletnya jauh di atas \qty{2.04}{ppm}: bukan senyawa yang dicari.

\textbf{12.} Empat sinyal (singlet jauh di kiri, triplet, sinyal 6
garis, triplet): bukan senyawa yang dicari, yang memiliki tiga sinyal.

\textbf{13.} Etil etanoat, \ce{CH3-CO-O-CH2-CH3}: singlet
\qty{2.04}{ppm} = \ce{CH3-C=O}; kuartet \qty{4.12}{ppm} = \ce{O-CH2};
triplet \qty{1.26}{ppm} = \ce{CH3} di ujung.

\textbf{14.} Singlet 3, kuartet 2, triplet 3.

\textbf{15.} $72 / 8 = \qty{9}{mm}$ per \omterm{def:g9:inside-the-atom:proton}{proton}.

\textbf{16.} Singlet $3 \times 9 = \qty{27}{mm}$; triplet
\qty{27}{mm}.

\textbf{17.} $18 + 27 + 27 = \qty{72}{mm}$.

\textbf{18.} Ya: \omterm{def:g11:functional-groups:carboxylic-acid}{ester-ester} kecil dikenal karena baunya yang harum
seperti buah.

\textbf{19.} Anak tangga kuartet itu \qty{18}{mm} dari \qty{72}{mm}.
\end{solution}
